\section{System overview \& layering}\label{sec:overview}

QtRocket is organized as three executables over one Qt-free core. \code{qtrocket} is the
Qt6 Widgets desktop application; \code{qtrocket-cli} is a headless, line-oriented
read-eval-print loop (REPL) over the same simulation engine; \code{qtrocket-visualizer} is
a standalone OpenGL viewer for saved designs. All three link a single static library,
\code{qtrocket\_core}, which contains exactly one class --- the \code{QtRocket}
application controller \src{CMakeLists.txt}{224} --- and pulls in the three layer
libraries beneath it: \code{model}, \code{sim}, and \code{utils}
\src{CMakeLists.txt}{229}. \Cref{fig:overview-layering} shows the arrangement at the
pinned commit \code{254cd41}; this section walks it from the top down and closes with the
build system that enforces it.

\tikzfig{system-layering}{Three executables over one Qt-free core. Every link edge points
strictly downward: the front ends link \code{qtrocket\_core} --- the static library
holding the \code{QtRocket} controller \srccap{CMakeLists.txt}{224} --- the core links the
\code{model}, \code{sim}, and \code{utils} layer libraries
\srccap{CMakeLists.txt}{229}, and \code{utils} links only the vendored Eigen and libcurl
\srccap{utils/CMakeLists.txt}{14}. Qt appears only in the two graphical executables;
nothing at or below the controller can name a Qt type.}{fig:overview-layering}

\subsection{Three executables, one core}\label{sec:overview-executables}

\begin{description}
\item[\code{qtrocket} (GUI).] The entry point sets the \code{Logger} to INFO
  \src{main.cpp}{16}, obtains the controller singleton \src{main.cpp}{20}, and blocks in
  \cppinline|gui::run(qtrocket, argc, argv)| \src{main.cpp}{25}. \code{gui::run} exists
  precisely so the controller stays Qt-free: it owns \code{QApplication}, and everything
  Qt-dependent in this executable lives under \srcf{gui/} \src{gui/GuiRunner.h}{13}. The executable links
  \code{Qt6::Widgets}, \code{Qt6::PrintSupport}, and \code{qtrocket\_core}
  \src{CMakeLists.txt}{286}.
\item[\code{qtrocket-cli} (headless REPL).] The same engine with no display requirement.
  Its entry point silences the logger to ERROR because stdout doubles as a
  machine-readable \code{OK}/\code{ERR}/\code{CSV} protocol \src{cli/CliMain.cpp}{17},
  then obtains the same singleton \src{cli/CliMain.cpp}{19}. The REPL is compiled into its
  own \code{cli} static library so the test binaries drive exactly the object code users
  type into \src{CMakeLists.txt}{234}; the executable links that library and nothing else
  \src{CMakeLists.txt}{297}.
\item[\code{qtrocket-visualizer} (OpenGL viewer).] A self-contained Qt application under
  \srcf{visualizer/} with its own \cppinline|find_package(Qt6 ...)| for the OpenGL
  components \src{visualizer/CMakeLists.txt}{9}; it links \code{qtrocket\_core} for the
  shared design loader \src{CMakeLists.txt}{241} and renders saved \code{.qrd} designs
  (\cref{sec:persistence}).
\end{description}

All three front ends are examined in depth in \cref{sec:interfaces}. The core they share
is deliberately minimal: \code{qtrocket\_core} compiles \srcf{QtRocket.cpp} exactly once
--- ``Application controller shared by the GUI, the CLI, and the integration tests''
\src{CMakeLists.txt}{222} --- and every consumer links the archive rather than recompiling
the translation unit (\cref{lst:overview-core}). The integration-test binary is the fourth
consumer \src{tests/CMakeLists.txt}{19}, so the controller's singleton statics exist once
per process no matter which front end, or which test, is running.

\begin{listing}[htbp]
\begin{minted}{cmake}
# Application controller shared by the GUI, the CLI, and the
# integration tests.
add_library(qtrocket_core STATIC
   QtRocket.cpp
   QtRocket.h)
qtrocket_enable_coverage(qtrocket_core)

target_link_libraries(qtrocket_core PUBLIC
   model
   sim
   utils)
\end{minted}
\caption{The keystone of the link graph \srccap{CMakeLists.txt}{222}. The controller
compiles once into a static library; the three executables and the integration tests all
link it, and its \code{PUBLIC} dependencies hand every consumer the layer libraries
below.}\label{lst:overview-core}
\end{listing}

\subsection{The layering is the link graph}\label{sec:overview-layers}

The layering is not a documentation convention; it is the link graph itself, declared
target by target in CMake. \code{utils} sits at the bottom and links only the two vendored
third-party libraries, \code{libcurl} and \code{eigen} \src{utils/CMakeLists.txt}{14} ---
it supplies the math typedefs (\code{Vector3}, \code{Matrix3}, \code{Quaternion}), the
\code{Logger}, and the HTTP transport everything above shares. \code{sim} links only
\code{utils} \src{sim/CMakeLists.txt}{28}. \code{model} links \code{utils}, \code{sim},
\code{Boost::property\_tree}, and \code{jsoncpp\_static} \src{model/CMakeLists.txt}{44};
the build file states why the middle edge exists --- ``sim is a real dependency:
Propagatable/RocketModel use sim::Aero and sim/StateData.h''
\src{model/CMakeLists.txt}{41}. \code{model} and \code{sim} are therefore layered rather
than peers: the model consumes the simulation layer's value types, and no \srcf{sim/} code
consumes a model value type in return. The one upward reference is the bridge interface
itself --- \srcf{sim/Propagator.h} includes \srcf{model/Propagatable.h}
\src{sim/Propagator.h}{18} to hold the abstract rocket it integrates, an abstraction whose
own includes point strictly downward \src{model/Propagatable.h}{12} --- so \code{sim}
still links only \code{utils}, and the \code{model} archive supplies the concrete rocket
when a final executable links both. Include resolution everywhere uses one global
\cppinline|include_directories(${PROJECT_SOURCE_DIR})| \src{CMakeLists.txt}{147}, which is
what makes the \code{model/...}, \code{sim/...}, and \code{utils/...} path prefixes
resolve uniformly from any file in the tree.

\begin{designrule}[The layering invariant, stated where it binds]
``Bottom layer: utils must not depend on model/ or sim/'' \src{utils/CMakeLists.txt}{13}.
The comment sits directly above the \cppinline|target_link_libraries| call it constrains,
and it holds: every \code{\#include} under \srcf{utils/} resolves to the standard library,
curl, Eigen, or another \srcf{utils/} header. This single rule keeps the dependency graph
acyclic --- every layer above may share the bottom layer's typedefs and logging without
ever creating a back-edge.
\end{designrule}

The invariant holds upward as well: nothing below the controller reaches back up to it.
The simulation environment enters the model as an explicit parameter of the force query
\src{model/Propagatable.h}{29}, passed by the propagator at the call site
\src{sim/Propagator.cpp}{41}; there is no \cppinline|QtRocket::getInstance()| call
anywhere in \srcf{model/} or \srcf{sim/}, so the layers below the controller are usable
--- and are unit-tested --- without a controller existing at all.

\subsection{The \code{QtRocket} controller}\label{sec:overview-controller}

\code{QtRocket} is the composition root: a lazily created, mutex-guarded singleton that
lives for the whole process and owns, through \cppinline|std::shared_ptr|, one
\code{(RocketModel, Propagator)} pair \src{QtRocket.h}{79}, one \code{sim::Environment}
\src{QtRocket.h}{82}, and one \code{model::MotorModelDatabase} \src{QtRocket.h}{83}. Its
constructor wires the system in four moves \src{QtRocket.cpp}{39}: the environment, whose
default constructor supplies working gravity and atmosphere models \src{QtRocket.cpp}{43}
(the alternatives are the subject of \cref{sec:environment}); the rocket model, which
boots with a placeholder body so the composite mass is real from the first query
\src{model/RocketModel.cpp}{30}; the propagator, constructed around that model with the
shared environment injected in a single statement \src{QtRocket.cpp}{48}; and the motor
database \src{QtRocket.cpp}{51}. Because the propagator is constructed together with the
model it integrates --- and the one API that would rebuild the pair, \code{addRocket}, has
no callers anywhere in the tree \src{QtRocket.h}{41} --- the model and the loop can never
drift apart, and the environment injection has exactly one site.

The singleton machinery is small enough to state completely. Three statics ---
\code{instance}, \code{mtx}, and a plain \code{bool initialized} \src{QtRocket.cpp}{14}
--- implement double-checked initialization: \code{getInstance()} reads the flag with no
lock \src{QtRocket.cpp}{21}; on false, \code{init()} acquires the mutex, re-checks under
the lock so a racing second caller cannot construct twice \src{QtRocket.cpp}{31}, and only
after allocating the controller publishes \cppinline|initialized = true|
\src{QtRocket.cpp}{35}. The instance is never deleted; it lives until the process exits.

\begin{keyinsight}[The non-atomic fast-path read]
The first check reads \code{initialized} --- a plain \code{bool}, not
\cppinline|std::atomic<bool>| --- outside the mutex \src{QtRocket.cpp}{21}, while
\code{init()} writes it under the lock \src{QtRocket.cpp}{35}. Concurrent first calls from
two threads would be a data race under the C++ memory model, and on weakly-ordered
hardware a second thread could observe the flag before the construction it guards. The
race is latent and benign today because every entry point that reaches the controller
calls \code{getInstance()} on the main thread before any other thread exists \src{main.cpp}{20},
\src{cli/CliMain.cpp}{19} --- safe by usage, not by construction. The fix idiom already
exists one layer down: the \code{Logger} is a Meyers singleton, initialized exactly once
by the language itself \src{utils/Logger.cpp}{19}.
\end{keyinsight}

The public surface is deliberately small, and almost all of it delegates:
\code{setTimeStep} and \code{setIntegratorModel} forward to the propagator
\src{QtRocket.h}{35}; \code{setInitialState} and \code{getStates} to the model
\src{QtRocket.h}{68}, \src{QtRocket.h}{48}; \code{getTrajectoryStatistics}
\src{QtRocket.h}{55} and the typed \code{getTerminationReason} \src{QtRocket.h}{62} return
whole-flight results; plain accessors expose the environment, the rocket, and the motor
database \src{QtRocket.h}{34}. The one method with a body worth reading is
\code{launchRocket()} \src{QtRocket.h}{43}, reproduced verbatim in
\cref{lst:overview-launch}: a flight is exactly four calls.

\begin{listing}[htbp]
\begin{minted}{cpp}
void QtRocket::launchRocket()
{
   // initialize the propagator
   rocket.first->clearStates();
   rocket.second->setCurrentTime(0.0);

   // start the rocket motor
   rocket.first->launch();

   // run the propagator until it terminates
   rocket.second->runUntilTerminate();
}
\end{minted}
\caption{The entire launch path in the controller \srccap{QtRocket.cpp}{54}: clear the
recorded history, rewind simulation time, ignite the model, run the loop to termination.
Everything else --- staging launch parameters, selecting models, exporting results ---
lives in the front ends.}\label{lst:overview-launch}
\end{listing}

\Cref{fig:overview-launch} places those four calls in the full request flow. In order:

\begin{enumerate}
\item \textbf{Clear history.} \cppinline|clearStates()| empties the
  $(t, \mathtt{StateData})$ vector stored on the model \src{model/Propagatable.h}{51}. Any
  element reference still held from a previous \code{getStates()} is invalidated here, which is
  why both front ends re-fetch the history after every launch
  \src{gui/AnalysisWindow.cpp}{39}, \src{cli/Repl.cpp}{690}.
\item \textbf{Rewind time.} \cppinline|setCurrentTime(0.0)| \src{sim/Propagator.h}{67}.
  Every flight starts at $t = 0$, which is also the ignition convention --- the simulation
  clock and the thrust-curve clock agree by construction.
\item \textbf{Ignite.} \cppinline|launch()| copies the staged initial state into the live
  state and arms the motor's thrust curve with \cppinline|startMotor(0.0)|
  \src{model/RocketModel.cpp}{105}. With no motor this is a no-op; the propagator's
  no-liftoff guard then ends the run with a typed verdict rather than a hang
  (\cref{sec:simcore}).
\item \textbf{Propagate.} \cppinline|runUntilTerminate()| integrates until the model's own
  \code{terminateCondition()} or one of the loop's guards fires \src{sim/Propagator.cpp}{58},
  resetting the whole-flight statistics at the top of the run \src{sim/Propagator.cpp}{65}.
  At every integrator stage the loop evaluates
  $\dot{\vec v} = \mathtt{getForces}(t, \vec x, \vec v, \mathtt{env}) /
  \mathtt{getMass}(t)$ \src{sim/Propagator.cpp}{41}; per recorded step it snapshots mass
  properties and appends to the history \src{sim/Propagator.cpp}{132}.
\end{enumerate}

The whole sequence is synchronous on the calling thread --- no worker thread, no lock
around the history, and nothing reads it concurrently (the consequences for the GUI are
taken up in \cref{sec:interfaces}). The canonical drive sequence ---
\code{setInitialState} $\to$ \code{setIntegratorModel} $\to$ \code{setTimeStep} $\to$
\code{launchRocket} $\to$ \code{getStates} --- appears verbatim in the integration
harness \src{tests/PhysicsIntegrationTests.cpp}{90}; the REPL issues the same calls, the
integrator and timestep staged by their own commands before its launch handler runs
\src{cli/Repl.cpp}{687}.

\tikzfig{seq-launch}{The launch sequence. A front end stages the flight through the
controller's delegating setters, then calls \code{launchRocket()}, whose entire body is
the four calls of \cref{lst:overview-launch} \srccap{QtRocket.cpp}{54}. Inside
\code{runUntilTerminate()} \srccap{sim/Propagator.cpp}{58} the rocket is visible only
through the \code{Propagatable} interface: force and mass at every integrator stage, a
mass-property snapshot per recorded step, and the termination predicate once per step.
Results flow back through const-reference accessors after the run
completes.}{fig:overview-launch}

\subsection{The \code{Propagatable} bridge}\label{sec:overview-bridge}

The seam between the two halves of the core is the abstract class
\code{model::Propagatable} \src{model/Propagatable.h}{23}: the only view of the rocket the
simulation loop ever sees. \code{RocketModel} is its only production implementation; the
\code{Propagator} holds it as \cppinline|std::shared_ptr<model::Propagatable>|
\src{sim/Propagator.h}{105} and interrogates it through six pure virtuals, reproduced
verbatim in \cref{lst:overview-propagatable}.

\begin{physbox}[What a simulation loop needs from a rocket]
A point-mass flight simulator integrates Newton's second law: given position and velocity,
it needs the net force $\vec F(t, \vec x, \vec v)$ and the current mass $m(t)$ --- mass is
a function of time because propellant burns away. A rigid-body simulator with all six
degrees of freedom (6-DOF: three translational plus three rotational) additionally needs
the net torque and the inertia tensor for the rotational equations. That is the entire
shape of the interface below: two questions for the translational problem, two for the
rotational one, plus a snapshot hook and an end-of-flight predicate.
\end{physbox}

\begin{listing}[htbp]
\begin{minted}{cpp}
   virtual Vector3 getForces(double t, const Vector3& position, const Vector3& velocity, sim::Environment& environment) = 0;
   virtual Vector3 getTorques(double t) = 0;

   virtual double getMass(double t) = 0;
   virtual Matrix3 getCompositeInertiaTensor(double t) = 0;

   /// Fill st.mass/cg/inertia with the body's composite mass properties at @p t. Called by the
   /// Propagator once per recorded step (before appendState).
   virtual void writeMassProperties(double t, StateData& st) = 0;

   virtual bool terminateCondition(double t) = 0;
\end{minted}
\caption{The model$\leftrightarrow$sim bridge, complete \srccap{model/Propagatable.h}{29}.
The integrator needs forces, torques, and time-varying mass properties --- and nothing
else about the rocket. \code{terminateCondition} lets the model, not the loop, decide when
a flight is over.}\label{lst:overview-propagatable}
\end{listing}

The bridge is more than a pure interface: it is also the state store. The base class
carries the current, initial, and recorded states, so the flight history lives on the
model side of the seam --- the propagator commits each accepted step through
\code{setCurrentState} \src{sim/Propagator.cpp}{124} and appends to the history exposed by
\code{getStates()} \src{model/Propagatable.h}{47}, which is why the controller's
\code{getStates()} delegates to the rocket rather than to the propagator
\src{QtRocket.h}{48}. The minimal contract is pinned by the propagator test suite: a
conforming mock implements only the six pure virtuals and inherits state storage, history,
and statistics from the base class \src{tests/PropagatorTests.cpp}{32}.

Of the six virtuals, four carry the current flight and two are dormant --- implemented,
with zero production consumers. The propagator's ordinary-differential-equation
(ODE) lambda consumes \code{getForces} and \code{getMass} at every integrator stage
\src{sim/Propagator.cpp}{41}; \code{writeMassProperties} runs once per recorded step
\src{sim/Propagator.cpp}{132}; \code{terminateCondition} is tested once per step after the
state is committed \src{sim/Propagator.cpp}{137}. \code{getTorques} is implemented as
constant zero \src{model/RocketModel.cpp}{94} and called by no simulation code --- the
comment beside the propagator's commented-out orientation integrator names it as the
future 6-DOF consumer \src{sim/Propagator.h}{102} --- and
\code{getCompositeInertiaTensor} likewise has no production caller; the per-step snapshot
reaches the part tree directly rather than through the virtual
\src{model/RocketModel.cpp}{59}. Two protected members are dead outright: \code{aeroData}
is default-constructed and never read \src{model/Propagatable.h}{61} --- the
\code{sim::Aero} alias survives specifically so it keeps compiling \src{sim/Aero.h}{63}
--- and \code{nextState} is never touched; the propagator uses its own local of the same
name \src{model/Propagatable.h}{65}, \src{sim/Propagator.cpp}{90}. The loop that drives
the live half is the subject of \cref{sec:simcore}; what the dormant halves are waiting
for is taken up in \cref{sec:aero}; \cref{sec:incomplete} carries the consolidated
inventory.

\subsection{Build system in brief}\label{sec:overview-build}

The build is one CMake project --- \cppinline|cmake_minimum_required(VERSION 3.25)|
\src{CMakeLists.txt}{1}, project \code{VERSION 0.1} \src{CMakeLists.txt}{3} --- pinned to
C++23 \src{CMakeLists.txt}{5} with \code{compile\_commands.json} exported for tooling
\src{CMakeLists.txt}{7}. Generator, compiler, and build type are all pinned by
\srcf{CMakePresets.json}: Ninja everywhere \src{CMakePresets.json}{7}, eight configure
presets (\code{debug}/\code{release} for the default compiler, Clang, and GCC; a
release-only MSVC preset \src{CMakePresets.json}{70}; and \code{coverage-clang}, which
redirects to a separate \code{build-coverage/} tree \src{CMakePresets.json}{35}), each
with matching build and test presets. The presets are the single sanctioned entry point
--- the command line, the editor tasks, and continuous integration (CI) all configure
through them --- a pinning adopted after a generator mismatch broke the build directory
\src{TODO.md}{24}.

Dependency policy is strict vendoring (\cref{tab:overview-deps}): every dependency except
Qt6 is fetched by \code{FetchContent} at an exact git tag and compiled from source at
first configure. A directory-scoped \cppinline|set(BUILD_SHARED_LIBS OFF)|
\src{CMakeLists.txt}{46} makes every in-tree library a static archive as well, so
\cref{fig:overview-layering} is a graph of static libraries resolved at link time, not a
runtime loading scheme.

\begin{table}[htbp]
\centering\small
\begin{tabularx}{\linewidth}{@{}ll>{\raggedright\arraybackslash}Xl@{}}
\toprule
Dependency & Pinned tag & Build shape and role & \\
\midrule
GoogleTest & \code{v1.17.0} & test framework for the six test binaries & \src{CMakeLists.txt}{17} \\
jsoncpp & \code{1.9.7} & static-only, tests off; thrustcurve.org JSON (\cref{sec:motors}) & \src{CMakeLists.txt}{32} \\
curl & \code{curl-8\_20\_0} & static, HTTP-only; the sole network transport, under \code{utils} & \src{CMakeLists.txt}{44} \\
Eigen & \code{5.0.1} & vectors, matrices, quaternions beneath \code{utils} & \src{CMakeLists.txt}{72} \\
Boost & \code{boost-1.91.0-1} & restricted to \code{property\_tree}; XML persistence (\cref{sec:persistence}) & \src{CMakeLists.txt}{79} \\
Qt6 & system & Widgets/PrintSupport for the GUI; OpenGL modules for the visualizer & \src{CMakeLists.txt}{140} \\
\bottomrule
\end{tabularx}
\caption{Dependency vendoring at the pinned commit. Everything except Qt6 is fetched at an
exact tag via \code{FetchContent}, so the dependency closure is a function of the
repository, not of the build host. The first configure downloads and compiles all of it
--- slow once, hermetic thereafter.}\label{tab:overview-deps}
\end{table}

\begin{designrule}[Hermetic dependency closure]
The curl block is the most deliberate piece of the vendoring. curl is built static and
HTTP-only \src{CMakeLists.txt}{46}, and every optional third-party backend (libssh2,
libssh, nghttp2, libidn2, brotli, zstd) is explicitly forced off, with the rationale recorded in
the tree: several of those options default to ON/AUTO, so curl links whatever it
auto-detects on the build host and the final link then fails on machines that lack those
libraries --- ``Forcing them off keeps the build identical on every platform''
\src{CMakeLists.txt}{56}. The payoff is stated where it pays: the FreeBSD CI job builds in
a bare virtual machine: ``The hermetic curl config means OpenSSL + zlib (both in
the FreeBSD base system) are the only external libs curl needs''
\src{.github/workflows/cmake-multi-platform.yml}{99}.
\end{designrule}

Warning hygiene is tree-wide and strict: \code{-Wall -Wextra -Wpedantic -Werror} on GCC
and Clang \src{CMakeLists.txt}{97}, \code{/W4 /WX} on MSVC \src{CMakeLists.txt}{88}. The
placement of those flags is itself the policy: they are added after the five
\code{FetchContent} blocks but before the project's own subdirectories, so vendored code
compiles with its own flags while every project target must be warning-clean; the one
vendored source inside a project target, \srcf{gui/qcustomplot.cpp}, is exempted
per-source with \code{-w} \src{CMakeLists.txt}{154}. The rest of the infrastructure ---
the six GoogleTest binaries and their dual ctest registration, the light/heavy test split,
the four-entry CI matrix plus the FreeBSD virtual-machine job, and the llvm-cov
coverage pipeline --- is the subject of \cref{sec:testing}.

\paragraph{Logging and API documentation.} Diagnostics everywhere in the tree flow through
the \code{utils::Logger} Meyers singleton \src{utils/Logger.cpp}{18}, which writes every
message to two sinks --- stdout and a \code{log.txt} opened in the process working
directory \src{utils/Logger.cpp}{25} --- at five cumulative levels up to \code{PERF\_}
\src{utils/Logger.h}{28}. Emission is mutex-guarded \src{utils/Logger.cpp}{35}; the level
setter is not \src{utils/Logger.cpp}{77}, which is safe under the current call pattern
(each front end pins the level once at startup, \cref{sec:interfaces}). A root
\srcf{Doxyfile} generates API documentation from the tree's \cppinline|///| comments.

\paragraph{Getting a build.} For the reader who wants the executables before the theory:

\begin{minted}{text}
cmake --preset debug-clang                          # configure (Ninja, into build/)
cmake --build --preset debug-clang                  # compile everything
ctest --test-dir build -R 'qtrocket_*' -LE heavy    # the fast test loop
\end{minted}

\noindent The first configure downloads and compiles the vendored closure of
\cref{tab:overview-deps}, so it is slow once and incremental thereafter. The executables
land at \code{build/qtrocket} and \code{build/qtrocket-cli}, with the viewer under
\code{build/visualizer/}; the test binaries and their ctest registration are mapped in
\cref{sec:testing}.

With the map established, the remaining sections descend into the boxes: the part tree
(\cref{sec:parts}) and placement (\cref{sec:placement}), motors (\cref{sec:motors}), the
simulation loop (\cref{sec:simcore}) and its environment models (\cref{sec:environment}),
aerodynamics (\cref{sec:aero}), the front ends (\cref{sec:interfaces}), and persistence
(\cref{sec:persistence}).
